% Copyright (c) 2026 Eugene Joseph M. Ragasa.
% SPDX-License-Identifier: Apache-2.0
% Author: Eugene Joseph M. Ragasa
% Formal proof completed: 2026-10-19
\documentclass[11pt]{article}
\usepackage[T1]{fontenc}
\usepackage{amsmath,amssymb,booktabs,tabularx}
\usepackage[margin=1in]{geometry}
\title{Proof Dossier: Coefficients of the Truncation Residual}
\author{Eugene Joseph M. Ragasa}
\date{Formal proof completed: 2026-10-19}
\begin{document}\maketitle
\noindent\begin{tabularx}{\textwidth}{@{}lX@{}}\toprule
Statement ID & \texttt{BG-FOURIER-THM-019} \\
Status & proved; no admissions \\
Lean declaration & \texttt{ProjectKoios.BandGauge.Fourier.}\newline\texttt{finiteFourierCoefficient\_truncationResidual} \\
Owning module & \texttt{ProjectKoios/BandGauge/Fourier/FiniteParseval.lean} \\
Axioms & \texttt{propext}, \texttt{Classical.choice}, \texttt{Quot.sound} \\
\bottomrule\end{tabularx}
\section{Statement and argument}
Every coefficient of full reconstruction minus support-truncated reconstruction is the original coefficient minus its truncated value. Use coefficient linearity and exact coefficient recovery in both terms.
\section{Boundary}
The result assumes the complete declared finite transform and does not certify an external mode inventory. Lean checked the exact statement; the dossier and proof were AI-assisted.
\end{document}
